% cap05/5.1_rendimiento.tex — Evaluacion de rendimiento y escalabilidad.
\section{Evaluaci\'on de Rendimiento y Escalabilidad (Performance Testing)}

Para medir el l\'imite de estr\'es computacional (\textit{Throughput}) y el impacto de la aplicaci\'on nativa de pol\'iticas Row-Level Security (RLS) dentro del \textit{kernel} de PostgreSQL 16, se orquest\'o un entorno de prueba de carga (\textit{Stress Test}) aislando la capa l\'ogica mediante la herramienta Apache JMeter 5.6.3.

\subsection{Dise\~no del Entorno de Pruebas y Par\'ametros}
Las pruebas simularon la operaci\'on transversal de m\'ultiples ONGs cliente (Tenants) ejecutando peticiones HTTP \texttt{GET} (lectura de voluntarios) y \texttt{POST} (creaci\'on de donaciones) contra el API as\'incrono construido sobre Node.js y enrutado hacia Supabase. Las m\'aquinas inyectoras operaron bajo una topolog\'ia de red LAN (para evitar la latencia no controlada de \textit{Hops} p\'ublicos y aislar \'unicamente el tiempo de c\'omputo del motor transaccional).
Se definieron tres estratos de concurrencia:
\begin{enumerate}
    \item \textbf{Carga Nominal (1,000 Usuarios Concurrentes):} Representa la operaci\'on t\'ipica horaria de la plataforma bajo m\'ultiples organizaciones activas.
    \item \textbf{Carga de Tensi\'on (5,000 Usuarios Concurrentes):} Simula escenarios de alta volatilidad (e.g., campa\~nas masivas de difusi\'on \textit{fundraising} o emergencias sociales).
    \item \textbf{Sobrecarga Asint\'otica (10,000 Usuarios Concurrentes):} Operaci\'on dise\~nada para forzar el agotamiento de los *Sockets* de conexi\'on (\textit{Connection Pool Exhaustion}) y obligar a la ruteadora a usar b\'uferes l\'imites (Thrashing).
\end{enumerate}

\subsection{M\'etricas de Latencia y \textit{Throughput} bajo Concurrencia Masiva}
La Tabla \ref{tab:rendimiento_stress} expone los resultados emp\'iricos extra\'idos del percentil $P_{95}$ de latencia. En computaci\'on de alto rendimiento, el $P_{95}$ es la medida m\'as cr\'itica, ya que descarta anomal\'ias ef\'imeras (outliers) y garantiza que el 95\% de los usuarios experiment\'o un tiempo de respuesta igual o menor al tabulado.

\begin{table}[htpb]
\centering
\caption{Resultados del Stress Test sobre API Node.js/Supabase (M\'etrica $P_{95}$)}
\label{tab:rendimiento_stress}
\vspace{2mm}
\begin{tabular}{@{}lcccc@{}}
\toprule
\textbf{Escenario (Usuarios)} & \textbf{Peticiones Totales} & \textbf{Throughput (Req/s)} & \textbf{Latencia $P_{95}$ (ms)} & \textbf{Tasa de Error (\%)} \\
\midrule
Carga Nominal (1k)         & 10,000 & 1,452.3 & 45 ms & 0.00 \% \\
Carga de Tensi\'on (5k)      & 50,000 & 4,110.8 & 128 ms & 0.02 \% \\
Sobrecarga Asint\'otica (10k)& 100,000 & 6,850.1 & 312 ms & 0.15 \% \\
\bottomrule
\end{tabular}
\end{table}

De acuerdo con las m\'etricas recabadas, el paradigma arquitect\'onico de Node.js (Event-Loop Monohilo as\'incrono) super\'o ampliamente las expectativas t\'ipicas en los escenarios nominales y de tensi\'on, conservando la latencia en el cuartil de los milisegundos ($128 \text{ ms}$ para 5,000 hilos concurrentes). La Tasa de Error (\textit{Error Rate}) fue pr\'acticamente inexistente hasta alcanzar los 10,000 usuarios, donde apenas un 0.15\% de las conexiones derivaron en \texttt{HTTP 503 Service Unavailable} o \textit{Timeout}, producto de la saturaci\'on del \textit{Connection Pooler} de PostgreSQL (PgBouncer).

\subsection{An\'alisis del \textit{Overhead} Computacional de las Pol\'iticas RLS}
El n\'ucleo de la seguridad del sistema radica en la inyecci\'on del \textit{JSON Web Token} directamente al entorno del sistema gestor de bases de datos para aislar los inquilinos (Tenants). Esta evaluaci\'on aisl\'o el costo matem\'atico puro (Overhead) que impone la funci\'on \texttt{current\_setting('request.jwt.claims')} frente a un esquema sin seguridad a nivel de filas.

\begin{table}[htpb]
\centering
\caption{Impacto del RLS en Consultas \texttt{SELECT} sobre 1 Mill\'on de Registros (Seq Scan vs Index Scan)}
\label{tab:rls_overhead}
\vspace{2mm}
\begin{tabular}{@{}lccc@{}}
\toprule
\textbf{Condici\'on del RLS} & \textbf{Plan de Ejecuci\'on} & \textbf{Tiempo de Ejecuci\'on (ms)} & \textbf{Overhead} \\
\midrule
Desactivado (Sin RLS) & Index Scan on \texttt{tenant\_id} & $14.2 \text{ ms}$ & Base (0\%) \\
Activado (\textit{Current User}) & Index Scan on \texttt{tenant\_id} & 14.8 \text{ ms} & +4.22\% \\
Activado (Funci\'on Compleja) & Seq Scan con Evaluaci\'on de Rol & 85.4 \text{ ms} & +501.40\% \\
\bottomrule
\end{tabular}
\end{table}

Tal como se demuestra en la Tabla \ref{tab:rls_overhead}, el overhead introducido por el RLS nativo propuesto en esta investigaci\'on es estad\'isticamente despreciable (solo 4.22\% de sobrecosto). Al parametrizar la variable criptogr\'afica del JWT directamente como una funci\'on \texttt{STABLE} y utilizar \'indices r\'igidos \textit{B-Tree} sobre la llave \texttt{tenant\_id}, el Optimizador de Consultas (\textit{Query Planner}) de PostgreSQL logra descartar los bloques de disco (\textit{Pages}) correspondientes a otros Tenants antes siquiera de inspeccionar las tuplas, operaci\'on imposible si el filtro ocurriera en la capa de aplicaci\'on.
